% Замените поясняющий текст собственными результатами.
% В шаблоне нет готового решения и обязательных внешних изображений.
\reportelement{Введение}

Цель работы - переносим запуск приложения, PostgreSQL и Redis в Docker Compose и пробуем масштабирование. В рамках лабораторной работы весь стек сервисов описывается в едином файле docker-compose.yml, поднимается одной командой и проверяется взаимодействие компонентов через пользовательскую сеть, именованные тома и механизм depends\_on с healthcheck. Дополнительно реализуется кэширование ответов внешнего API погоды в Redis и демонстрируется горизонтальное масштабирование stateless-сервиса приложения.

\begin{table}[H]
  \caption{Персональные параметры}\label{tab:parameters}
  \centering
  \begin{tabularx}{\textwidth}{|p{0.46\textwidth}|X|}
    \hline
    Параметр & Значение\\ \hline
    Номер ИСУ & \StudentISU\\ \hline
    Персональный порт & 5455, 8025, 6395\\ \hline
    Персональный префикс & manuk\\ \hline
    Секретный ключ & secret25manuk\\ \hline
    Имя баз данных & manuk-postgres\\ \hline
  \end{tabularx}
\end{table}


Результат хранится в ветке \texttt{lab-\LabNumber} своего форка:
\url{\RepositoryURL}.

\chapter{Ход выполнения}
\section{Подготовка проекта для Docker Compose}

Создаем файл env с персональными переменными, данный файл не попадет в гит, так как во второй работе добавили этот тип файлов в .gitignore для сохранности данных. 

\begin{code}
APP_OWNER=manuk
APP_SECRET=secret25manuk
APP_PORT=8025
POSTGRES_USER=user_25
POSTGRES_PASSWORD=pgpass_25
POSTGRES_DB=assistant_25
PG_PORT=5455
REDIS_PORT=6395
CACHE_TTL=120
\end{code}

Создали файл docker-compose.yml в корне проекта. 

\begin{code}
services:
  db:
    image: postgres:15
    container_name: manuk-postgres
    environment:
      POSTGRES_USER: ${POSTGRES_USER}
      POSTGRES_PASSWORD: ${POSTGRES_PASSWORD}
      POSTGRES_DB: ${POSTGRES_DB}
    ports:
      - "${PG_PORT}:5432"
    volumes:
      - manuk-pgdata:/var/lib/postgresql/data
    networks:
      - manuk-backend
    healthcheck:
      test: ["CMD-SHELL", "pg_isready -U ${POSTGRES_USER}"]
      interval: 30s
      timeout: 5s
      retries: 3

  app:
    build: .
    container_name: manuk-assistant
    ports:
      - "${APP_PORT}:8025"  
    environment:
      APP_OWNER: ${APP_OWNER}
      APP_SECRET: ${APP_SECRET}
      DATABASE_URL: postgresql://${POSTGRES_USER}:${POSTGRES_PASSWORD}@db:5432/${POSTGRES_DB}
      CACHE_TTL: ${CACHE_TTL}
    depends_on:
      db:
        condition: service_healthy
    networks:
      - manuk-backend

volumes:
  manuk-pgdata:

networks:
  manuk-backend:
    driver: bridge
\end{code}

Запускаем контейнеры и проверяем работоспособность приложения. (рисунок \ref{fig:1})

\begin{code}
docker compose up -d --build
docker compose ps
docker compose logs app

curl http://localhost:8025/
curl http://localhost:8025/health
\end{code}

\reportscreenshot{pic/1.png}{Терминальный вывод}{fig:1}



\reportelement{Заключение}

Цель работы достигнута: запуск приложения, PostgreSQL и Redis перенесён в Docker Compose и выполнено масштабирование. Стек из трёх сервисов успешно оркестрируется, приложение стартует только после готовности базы и Redis, кэш погоды работает корректно (первый запрос идёт в API, последующие возвращаются из Redis), а при масштабировании до трёх экземпляров app все они используют общий кэш. Таким образом, показаны основные возможности Docker Compose как инструмента локальной оркестрации контейнеров.

% Оставьте только использованные источники и цитируйте их в тексте.
\begin{thebibliography}{9}
  \bibitem{docker-docs}
  Docker Documentation: документация Docker.
  URL: \url{https://docs.docker.com/}
  (дата обращения: укажите свою дату).
\end{thebibliography}


